\documentclass[runningheads]{llncs}
\usepackage[utf8]{inputenc}
\usepackage{graphicx}
\usepackage{parskip}
\usepackage{hyperref} 
\usepackage{float}


\begin{document}

\title{mySCDN - FlashCacheCDN}
\author{Iancu Stefan-Teodor}
\institute{Faculty of Computer Science, Alexandru Ioan Cuza University of Iași, Romania}
\maketitle

\section{Introduction}
\subsection{General Vision}
The proposal involves the creation of a CDN (Content Delivery Network) system aimed at delivering a service or transferring data as quickly as possible. Communication latency will be minimized, providing the user with a good application experience. Additionally, the goal of this project (acquiring hardware, deployment services, etc., enabling the entire system) is to generate minimal costs while delivering quality software services. The project solution proposes distributing multiple nodes in different geographical locations that will serve clients optimally.

\subsection{Objectives}
The project has the following objectives:
\begin{itemize}
    \item Prioritizing the customer experience.
    \item Allocating the most convenient node for each client.
    \item Avoiding the risk of compromising the entire service and application by always having a backup option (if a node crashes, another node will provide the service with latency acceptable to the client).
    \item Managing concurrency as reliably as possible.
    \item Software optimization to minimize CPU usage.
    \item Gradual optimization of the caching system and resources generated by concurrency.
    \item Avoiding unforeseen cases and handling errors effectively.
\end{itemize}

\section{Applied Technologies}
\begin{itemize}
    \item At the transport level, the project will use the TCP protocol throughout the system. Although it is more costly in terms of packet size and transmission time for a valid message exchange, it ensures that no packet will be lost or corrupted (internally). The project cannot afford, for example, to mistakenly allocate a node in America for a client in Madagascar due to the corruption of 3 bytes in a packet containing the IPv4 address of the node.
    
    \item To store information about the available cache on each node, the project will use the JSON format, allowing content to be deserialized into objects stored in memory. Server and client configuration files will be of type \texttt{.ini} because they are lightweight (as is the API) and easy to parse. Consequently, two open-source APIs for the C++ language were used, referenced in the bibliography notes.

    \item The architecture is designed for Unix-based systems. For network communication, standard C/C++ libraries for Unix sockets were utilized.

    \item Regarding concurrency, servers will implement multithreading, multiprocessing, and multiplexing depending on the case. For this, the standard libraries \texttt{<thread>} in C++ and \texttt{<unistd.h>} in C were used for management.
    
    \item The application logic was implemented using data structures provided by the standard C++ library and the OOP paradigm.
    
    \item CMake and bash scripts were used for easy builds and dependency installation.

\end{itemize}

\section{Application Structure}
The application architecture will be based on 4 entities: 
- Main server 
- Edge server
- EDNS server
- Client

A brief description of each:
\begin{itemize}

\item The main server will be responsible for storing all original resources, providing updates to edge servers, keeping track of cache and load on the edges, and transmitting it to the authoritative EDNS server.

\item Edge servers will cache certain resources to serve clients while transmitting load information to the main server.

\item The EDNS server will receive information about all edge servers and map them to make decisions based on cache, load, and geographic location. (it will also act as a load balancer)

\item The client will make a request for the desired resource to the EDNS server, receive an IP address in response, and connect to the edge server with the specified IP address.

\end{itemize}

In general terms, the architecture of this type of application can be very well described through an image:

\begin{figure} [H]
    \centering
    \includegraphics[width=1\linewidth]{with-CDN.jpeg}
    \caption{CDN over the world}
    \label{fig:enter-label}
\end{figure}

Now let's look at the next diagram (zoomed in):

\begin{figure}[H] 
    \centering
    \includegraphics[width=1\linewidth]{first_diagram.png}
    \caption{CDN Architecture}
    \label{fig:enter-label}
\end{figure}


\section{Implementation Aspects}

\subsection{Application-Level Protocol and Innovative Code Sections}

\textbf{Client-EDNS Server Protocol:}
\begin{itemize}
    \item The request packet will consist of: packet length + client IP + resource + null byte
    \item The response will be 4 bytes representing the edge server's IP address
\end{itemize}
\begin{figure}
    \centering
    \includegraphics[width=0.9\linewidth]{packet client.png}
    \caption{Example Client-EDNS server packet}
    \label{fig:enter-label}
\end{figure}

\textbf{Client-Edge Server Protocol:}
\begin{itemize}
    \item The request packet will consist of: packet length + resource + null byte
    \item The response will be packet length + response + null byte
    \item \textit{Similar to Fig.3}
\end{itemize}

\textbf{Edge Server-Main Server Protocol:}
\begin{itemize}
    \item The edge server will periodically transmit information about itself through a JSON file transfer or send a GET request to the server to request a resource in case of a cache MISS.
    \item The main server will send certain resources to the edge server along with a TTL (time to live) quantum.
    \item \textit{Similar to Fig.3} but file transfer will be split into byte chunks.
\end{itemize}

\textbf{Main Server-EDNS Server Protocol:}
\begin{itemize}
    \item The main server will transmit a JSON file with data about cache and load from edge servers.
    \item \textit{Similar to Fig.3} but file transfer will be split into byte chunks.
\end{itemize}

\begin{figure} [H]
    \centering
    \includegraphics[width=0.75\linewidth]{json.png}
    \caption{JSON data}
    \label{fig:enter-label}
\end{figure}

\begin{figure} [H]
    \centering
    \includegraphics[width=1\linewidth]{cache.png}
    \caption{Level 1 Deserialization}
    \label{fig:enter-label}
\end{figure}

\begin{figure} [H]
    \centering
    \includegraphics[width=1\linewidth]{edge_server.png}
    \caption{Level 2 Deserialization}
    \label{fig:enter-label}
\end{figure}

\subsection{Concurrency}

\begin{itemize}
    \item \textbf{Main server} - At startup, it will perform \textbf{pre-forking} to handle the N edge servers, while the root process will communicate with the EDNS server. This approach was chosen because the project cannot afford for the Main server to fail if one process crashes. Even though the forking process is costly, it is done during initialization and does not affect concurrent communication.

    \item \textbf{Edge server} - It will contain two processes. One will communicate with the Main server, and the other will create a \textbf{thread pool} of size \textit{k} (which will be gradually optimized for accuracy). Each thread runs the same function to serve clients, invoking the \texttt{accept()} primitive inside a \texttt{mutex.lock()}. This approach is used because creating a thread on the spot for high traffic volumes can cost a few milliseconds and put a load on the CPU, even though threads are lightweight. This implementation will also depend on whether the deployment service supports \textbf{multithreading}.

    \item \textbf{EDNS server} - It will have two threads with simultaneous access to a \texttt{Cache} variable allocated on the heap. One thread will modify the variable when receiving updates from the Main server, while the other will handle clients in a \textbf{multiplexed} manner and access the variable to decide on the most suitable edge server. Of course, \texttt{mutex.lock()} will be used to protect against \textit{undefined behavior} (e.g., one thread modifies while the other reads). This threading approach allows simultaneous access to \texttt{Cache}, and multiplexing is justified by reducing costs (CPU calls) and the lightweight nature of DNS server services.

    \item \textbf{Client} - Two processes. One will receive input and make requests to the EDNS server, then create a new process by calling \texttt{fork()} and redirect the request to the edge server, awaiting output. This approach allows a request not to be limited by the processing time of the previous (or other previous) requests.
\end{itemize}

In conclusion, the entire architecture of the developed CDN system can be encapsulated in the following detailed diagram (zoomed in):
\begin{figure}[H]
    \centering
    \includegraphics[width=1\linewidth]{second_diagram.png}
    \caption{Detailed Architecture}
    \label{fig:enter-label}
\end{figure}

\subsection{Real-world Use Case}

An aviation company wants to distribute its online booking service worldwide because it is popular and owns many airplanes. The company has two headquarters where servers are located, one in Europe and one in the USA. Data analysts hired by the company discovered that although clients traveling in Asia, Africa, or South America are very satisfied with the services, they do not generate much profit. Another study revealed that individuals who experience a web page load time of at least 5 seconds have a low probability of purchasing a flight ticket. Consequently, the company deploys its web service on edge servers strategically placed in geographic regions worldwide, provided by another company for a fee. In the following months, profits began to increase, and server load at headquarters was reduced because web traffic is now distributed optimally. Additionally, the risk of DoS and DDoS attacks was minimized.

\section{Potential Improvements}

\begin{itemize}
    \item Anti-DoS and DDoS protection, such as a proxy rate limiter for the EDNS server.
    \item In addition to TTL, implementing a caching algorithm on edge servers, such as LFU, LRU, etc.
    \item Adding a caching mechanism for lookups on the EDNS server.
    \item Developing an encryption tunnel for communication between servers, such as TLS.
\end{itemize}

\section{References}
\begin{itemize}
    \item \url{https://edu.info.uaic.ro/computer-networks/}
    \item \url{https://www.cloudflare.com/learning/cdn/what-is-a-cdn/}
    \item \url{https://aws.amazon.com/what-is/cdn/}
    \item \url{https://en.wikipedia.org/wiki/Multithreading_(computer_architecture)}
    \item \url{https://www.youtube.com/watch?v=pdPSLm629yk&t=454s}
    \item \url{https://en.wikipedia.org/wiki/Content_delivery_network}
    \item \url{https://en.cppreference.com/w/cpp/thread/thread}
    \item \url{https://en.cppreference.com/w/cpp/thread/mutex}
    \item \url{https://en.cppreference.com/w/cpp/container/unordered_set}
    \item \url{https://github.com/nlohmann/json}
    \item \url{https://github.com/dzilles/configparser}
\end{itemize}

\end{document}
